% ================= SEKCJA 6 =================
\section{Scenariusze}

W niniejszym rozdziale przedstawiono scenariusze wybranych przypadków użycia systemu \texttt{WorkFlow}. Scenariusze opisują przebieg najważniejszych operacji z punktu widzenia użytkownika lub komponentu technicznego, a także wskazują warunki początkowe, przebieg podstawowy oraz możliwe warianty alternatywne.

Opisane scenariusze odpowiadają aktualnemu zakresowi projektu i obejmują logowanie, zarządzanie profilem pracownika, rejestrację czasu pracy z czytnika RCP, obsługę certyfikatów, nieobecności, umów, dokumentów oraz przygotowanie danych rozliczeniowych.

\subsection{SC-01: Logowanie użytkownika}
\begin{itemize}
    \item \textbf{Numer:} SC-01
    \item \textbf{Nazwa:} Logowanie użytkownika
    \item \textbf{Aktor:} Użytkownik systemu
    \item \textbf{Stan wejścia:}
    \begin{itemize}
        \item Użytkownik posiada konto w systemie.
        \item Konto ma przypisany adres e-mail, hash hasła oraz rolę systemową.
        \item Konto jest aktywne i ma włączoną możliwość logowania.
    \end{itemize}
    \item \textbf{Przebieg:}
    \begin{enumerate}
        \item Użytkownik otwiera formularz logowania.
        \item Wprowadza adres e-mail oraz hasło.
        \item System przesyła dane logowania do backendu.
        \item Backend wyszukuje użytkownika po adresie e-mail.
        \item Backend weryfikuje poprawność hasła.
        \item Po poprawnej weryfikacji system generuje token JWT.
        \item Token zostaje zapisany w ciasteczku sesyjnym.
        \item Użytkownik zostaje przekierowany do panelu aplikacji.
    \end{enumerate}
    \item \textbf{Zakończenie lub przebieg alternatywny:}
    \begin{itemize}
        \item \textbf{Zakończenie pozytywne:} Użytkownik zostaje zalogowany i uzyskuje dostęp do funkcji zgodnych ze swoją rolą.
        \item \textbf{Błędne dane logowania:} System odrzuca próbę logowania i wyświetla komunikat o błędnych poświadczeniach.
        \item \textbf{Konto nieaktywne lub bez dostępu do logowania:} System nie przyznaje dostępu do aplikacji.
    \end{itemize}
\end{itemize}

\subsection{SC-02: Zarządzanie profilem pracownika}
\begin{itemize}
    \item \textbf{Numer:} SC-02
    \item \textbf{Nazwa:} Zarządzanie profilem pracownika
    \item \textbf{Aktor:} HR lub Administrator
    \item \textbf{Stan wejścia:}
    \begin{itemize}
        \item Użytkownik jest zalogowany.
        \item Użytkownik posiada rolę umożliwiającą zarządzanie danymi pracowników.
        \item W systemie istnieje zakład, do którego można przypisać pracownika.
    \end{itemize}
    \item \textbf{Przebieg:}
    \begin{enumerate}
        \item Użytkownik przechodzi do listy pracowników.
        \item System wyświetla listę pracowników zgodną z zakresem dostępu użytkownika.
        \item Użytkownik wyszukuje pracownika lub otwiera formularz dodania nowej osoby.
        \item Użytkownik wprowadza albo modyfikuje dane podstawowe pracownika.
        \item Użytkownik określa zakład, rolę, status aktywności oraz opcjonalny identyfikator karty RFID.
        \item System waliduje poprawność danych, między innymi unikalność adresu e-mail, numeru PESEL i identyfikatora karty RFID.
        \item System zapisuje dane w bazie.
        \item Użytkownik może przejść do profilu pracownika i zarządzać powiązanymi informacjami.
    \end{enumerate}
    \item \textbf{Zakończenie lub przebieg alternatywny:}
    \begin{itemize}
        \item \textbf{Zakończenie pozytywne:} Dane pracownika zostają zapisane lub zaktualizowane.
        \item \textbf{Błąd walidacji:} System wyświetla komunikat i wskazuje pola wymagające poprawy.
        \item \textbf{Brak uprawnień:} System blokuje operację, jeżeli użytkownik nie posiada odpowiedniej roli lub dostępu do zakładu.
    \end{itemize}
\end{itemize}

\subsection{SC-03: Rejestracja czasu pracy przez czytnik RCP}
\begin{itemize}
    \item \textbf{Numer:} SC-03
    \item \textbf{Nazwa:} Rejestracja czasu pracy przez czytnik RCP
    \item \textbf{Aktor:} Czytnik RCP
    \item \textbf{Stan wejścia:}
    \begin{itemize}
        \item W systemie istnieje aktywny czytnik RCP.
        \item Czytnik jest przypisany do zakładu i projektu.
        \item Pracownik posiada przypisany identyfikator karty RFID.
        \item Endpoint urządzenia jest zabezpieczony tokenem technicznym.
    \end{itemize}
    \item \textbf{Przebieg:}
    \begin{enumerate}
        \item Pracownik odbija kartę RFID na czytniku.
        \item Czytnik wysyła do backendu zdarzenie zawierające numer seryjny czytnika, identyfikator karty, typ zdarzenia oraz czas zdarzenia.
        \item Backend weryfikuje token urządzenia.
        \item System wyszukuje aktywny czytnik na podstawie numeru seryjnego.
        \item System wyszukuje aktywnego pracownika na podstawie identyfikatora karty RFID.
        \item System sprawdza, czy pracownik ma dostęp do zakładu powiązanego z czytnikiem.
        \item System zapisuje surowe zdarzenie w tabeli \texttt{TimeEvent}.
        \item Jeżeli zdarzenie ma typ \texttt{OUT}, system wyszukuje wcześniejsze zdarzenie \texttt{IN}.
        \item Na podstawie pary zdarzeń \texttt{IN}/\texttt{OUT} system tworzy wpis czasu pracy w tabeli \texttt{TimeEntry}.
    \end{enumerate}
    \item \textbf{Zakończenie lub przebieg alternatywny:}
    \begin{itemize}
        \item \textbf{Zakończenie pozytywne:} Zdarzenie zostaje zapisane, a dla zdarzenia \texttt{OUT} powstaje wpis czasu pracy.
        \item \textbf{Nieznany czytnik:} System odrzuca zdarzenie, jeżeli numer seryjny nie odpowiada aktywnemu czytnikowi.
        \item \textbf{Nieznana karta RFID:} System odrzuca zdarzenie, jeżeli nie znajdzie aktywnego pracownika.
        \item \textbf{Brak wcześniejszego wejścia:} System zapisuje błąd logiczny i nie tworzy wpisu czasu pracy dla zdarzenia \texttt{OUT}.
    \end{itemize}
\end{itemize}

\subsection{SC-04: Weryfikacja wpisów czasu pracy}
\begin{itemize}
    \item \textbf{Numer:} SC-04
    \item \textbf{Nazwa:} Weryfikacja wpisów czasu pracy
    \item \textbf{Aktor:} HR, Administrator lub użytkownik uprawniony do weryfikacji
    \item \textbf{Stan wejścia:}
    \begin{itemize}
        \item W systemie istnieją wpisy czasu pracy.
        \item Wpisy posiadają status, na przykład \texttt{PENDING}.
        \item Użytkownik posiada uprawnienia do przeglądania i weryfikowania wpisów czasu pracy.
    \end{itemize}
    \item \textbf{Przebieg:}
    \begin{enumerate}
        \item Użytkownik przechodzi do widoku czasu pracy.
        \item System wyświetla wpisy czasu pracy z wybranego zakresu.
        \item Użytkownik analizuje dane pracownika, projekt, czas rozpoczęcia, czas zakończenia oraz liczbę godzin.
        \item Użytkownik zatwierdza poprawny wpis albo oznacza go jako wymagający odrzucenia.
        \item System aktualizuje status wpisu czasu pracy.
        \item Zatwierdzone wpisy mogą zostać wykorzystane w module rozliczeń.
    \end{enumerate}
    \item \textbf{Zakończenie lub przebieg alternatywny:}
    \begin{itemize}
        \item \textbf{Zakończenie pozytywne:} Wpis uzyskuje status \texttt{APPROVED} i może zostać uwzględniony w rozliczeniach.
        \item \textbf{Odrzucenie wpisu:} Wpis uzyskuje status \texttt{REJECTED} i nie jest brany pod uwagę przy eksporcie payroll.
        \item \textbf{Brak uprawnień:} System blokuje zmianę statusu wpisu.
    \end{itemize}
\end{itemize}

\subsection{SC-05: Dodanie umowy pracownika}
\begin{itemize}
    \item \textbf{Numer:} SC-05
    \item \textbf{Nazwa:} Dodanie umowy pracownika
    \item \textbf{Aktor:} HR lub Administrator
    \item \textbf{Stan wejścia:}
    \begin{itemize}
        \item Użytkownik jest zalogowany.
        \item Użytkownik posiada dostęp do profilu pracownika.
        \item Pracownik istnieje w systemie.
    \end{itemize}
    \item \textbf{Przebieg:}
    \begin{enumerate}
        \item Użytkownik otwiera profil pracownika.
        \item Przechodzi do sekcji umów.
        \item Wybiera opcję dodania nowej umowy.
        \item Wprowadza typ umowy, kwotę wynagrodzenia, datę rozpoczęcia oraz opcjonalną datę zakończenia.
        \item Określa, czy umowa jest aktualna.
        \item System waliduje dane formularza.
        \item System zapisuje umowę i przypisuje ją do pracownika.
    \end{enumerate}
    \item \textbf{Zakończenie lub przebieg alternatywny:}
    \begin{itemize}
        \item \textbf{Zakończenie pozytywne:} Umowa zostaje zapisana w historii umów pracownika.
        \item \textbf{Brak wymaganych danych:} System wyświetla komunikat i nie zapisuje umowy.
        \item \textbf{Brak uprawnień:} System blokuje operację.
    \end{itemize}
\end{itemize}

\subsection{SC-06: Monitorowanie ważności certyfikatów}
\begin{itemize}
    \item \textbf{Numer:} SC-06
    \item \textbf{Nazwa:} Monitorowanie ważności certyfikatów
    \item \textbf{Aktor:} HR lub Administrator
    \item \textbf{Stan wejścia:}
    \begin{itemize}
        \item W systemie istnieje słownik certyfikatów.
        \item Pracownicy posiadają przypisane certyfikaty z datą wydania i datą ważności.
        \item Użytkownik posiada dostęp do modułu certyfikatów.
    \end{itemize}
    \item \textbf{Przebieg:}
    \begin{enumerate}
        \item Użytkownik przechodzi do widoku certyfikatów wygasających.
        \item Wybiera zakres czasu, na przykład 30, 90 lub 365 dni.
        \item System pobiera certyfikaty, których data ważności mieści się w wybranym zakresie.
        \item System prezentuje dane pracownika, typ certyfikatu, nazwę dokumentu oraz liczbę dni do wygaśnięcia.
        \item Użytkownik może przejść do profilu pracownika i uzupełnić lub zaktualizować dane certyfikatu.
    \end{enumerate}
    \item \textbf{Zakończenie lub przebieg alternatywny:}
    \begin{itemize}
        \item \textbf{Zakończenie pozytywne:} Użytkownik otrzymuje listę certyfikatów wymagających kontroli.
        \item \textbf{Brak wyników:} System wyświetla pustą listę dla wybranego zakresu.
        \item \textbf{Brak dostępu do zakładu:} System nie pokazuje danych spoza zakresu uprawnień użytkownika.
    \end{itemize}
\end{itemize}

\subsection{SC-07: Dodanie nieobecności pracownika}
\begin{itemize}
    \item \textbf{Numer:} SC-07
    \item \textbf{Nazwa:} Dodanie nieobecności pracownika
    \item \textbf{Aktor:} HR lub Administrator
    \item \textbf{Stan wejścia:}
    \begin{itemize}
        \item Użytkownik jest zalogowany.
        \item Pracownik istnieje w systemie.
        \item Użytkownik posiada dostęp do profilu pracownika.
    \end{itemize}
    \item \textbf{Przebieg:}
    \begin{enumerate}
        \item Użytkownik otwiera profil pracownika.
        \item Przechodzi do sekcji nieobecności.
        \item Wybiera opcję dodania nowej nieobecności.
        \item Określa typ nieobecności, datę rozpoczęcia i datę zakończenia.
        \item Opcjonalnie wskazuje dokument powiązany z nieobecnością.
        \item Określa, czy nieobecność jest zatwierdzona.
        \item System zapisuje nieobecność i przypisuje ją do pracownika.
    \end{enumerate}
    \item \textbf{Zakończenie lub przebieg alternatywny:}
    \begin{itemize}
        \item \textbf{Zakończenie pozytywne:} Nieobecność pojawia się w profilu pracownika.
        \item \textbf{Niepoprawny zakres dat:} System odrzuca zapis i wyświetla komunikat.
        \item \textbf{Brak uprawnień:} System blokuje operację.
    \end{itemize}
\end{itemize}

\subsection{SC-08: Dodanie dokumentu pracownika}
\begin{itemize}
    \item \textbf{Numer:} SC-08
    \item \textbf{Nazwa:} Dodanie dokumentu pracownika
    \item \textbf{Aktor:} HR lub Administrator
    \item \textbf{Stan wejścia:}
    \begin{itemize}
        \item Użytkownik posiada dostęp do profilu pracownika.
        \item Pracownik istnieje w systemie.
        \item Magazyn plików jest skonfigurowany.
    \end{itemize}
    \item \textbf{Przebieg:}
    \begin{enumerate}
        \item Użytkownik otwiera profil pracownika.
        \item Wybiera opcję dodania dokumentu.
        \item Wskazuje plik dokumentu lub dane dokumentu.
        \item System zapisuje plik w magazynie obiektowym albo zapisuje jego ścieżkę.
        \item System tworzy rekord dokumentu w bazie danych.
        \item Dokument zostaje powiązany z pracownikiem oraz opcjonalnie z umową lub certyfikatem.
    \end{enumerate}
    \item \textbf{Zakończenie lub przebieg alternatywny:}
    \begin{itemize}
        \item \textbf{Zakończenie pozytywne:} Dokument zostaje zapisany i widoczny w danych pracownika.
        \item \textbf{Błąd zapisu pliku:} System informuje użytkownika o problemie z magazynem plików.
        \item \textbf{Brak danych wymaganych:} System nie tworzy rekordu dokumentu.
    \end{itemize}
\end{itemize}

\subsection{SC-09: Przygotowanie eksportu rozliczeniowego}
\begin{itemize}
    \item \textbf{Numer:} SC-09
    \item \textbf{Nazwa:} Przygotowanie eksportu rozliczeniowego
    \item \textbf{Aktor:} Księgowość
    \item \textbf{Stan wejścia:}
    \begin{itemize}
        \item Użytkownik jest zalogowany i posiada rolę księgowości.
        \item W systemie istnieją zatwierdzone wpisy czasu pracy.
        \item Wpisy nie zostały wcześniej powiązane z eksportem payroll.
    \end{itemize}
    \item \textbf{Przebieg:}
    \begin{enumerate}
        \item Użytkownik przechodzi do modułu rozliczeń.
        \item Wybiera miesiąc i rok rozliczeniowy.
        \item System pobiera zatwierdzone wpisy czasu pracy z wybranego okresu.
        \item System tworzy rekord eksportu rozliczeniowego.
        \item System zapisuje liczbę pracowników objętych eksportem.
        \item System wiąże wybrane wpisy czasu pracy z eksportem.
        \item Użytkownik otrzymuje informację o przygotowaniu danych.
    \end{enumerate}
    \item \textbf{Zakończenie lub przebieg alternatywny:}
    \begin{itemize}
        \item \textbf{Zakończenie pozytywne:} Eksport zostaje utworzony, a powiązane wpisy czasu pracy są oznaczone jako rozliczone w ramach danej paczki.
        \item \textbf{Brak zatwierdzonych wpisów:} System nie tworzy eksportu i informuje użytkownika o braku danych.
        \item \textbf{Dane już wyeksportowane:} System pomija wpisy powiązane z wcześniejszym eksportem.
    \end{itemize}
\end{itemize}

\subsection{SC-10: Uruchomienie danych demonstracyjnych}
\begin{itemize}
    \item \textbf{Numer:} SC-10
    \item \textbf{Nazwa:} Uruchomienie danych demonstracyjnych
    \item \textbf{Aktor:} Administrator techniczny
    \item \textbf{Stan wejścia:}
    \begin{itemize}
        \item Środowisko Docker Compose jest uruchomione.
        \item Backend ma dostęp do bazy danych.
        \item W środowisku ustawiono zmienną potwierdzającą zgodę na wykonanie seeda demonstracyjnego.
    \end{itemize}
    \item \textbf{Przebieg:}
    \begin{enumerate}
        \item Administrator techniczny uruchamia komendę seeda demonstracyjnego.
        \item Skrypt sprawdza obecność wymaganej zmiennej środowiskowej.
        \item Skrypt czyści dane przeznaczone do prezentacji.
        \item Skrypt tworzy zakłady, użytkowników, pracowników, projekty, czytniki RCP, umowy, certyfikaty, dokumenty, nieobecności, zdarzenia czasu pracy i eksport rozliczeniowy.
        \item Administrator loguje się do aplikacji na konto testowe.
        \item System prezentuje komplet przykładowych danych.
    \end{enumerate}
    \item \textbf{Zakończenie lub przebieg alternatywny:}
    \begin{itemize}
        \item \textbf{Zakończenie pozytywne:} Baza danych zawiera komplet danych demonstracyjnych gotowych do prezentacji projektu.
        \item \textbf{Brak zmiennej zabezpieczającej:} Skrypt przerywa działanie i nie modyfikuje bazy.
        \item \textbf{Błąd połączenia z bazą:} Skrypt kończy się błędem i wyświetla komunikat diagnostyczny.
    \end{itemize}
\end{itemize}